\documentclass[aspectratio=169,11pt]{beamer}
\usetheme{Madrid}
\usecolortheme{dolphin}
\usepackage{fontspec}
\usepackage{xeCJK}
\usepackage{graphicx}
\usepackage{booktabs}
\setmainfont{Noto Sans CJK SC}
\setsansfont{Noto Sans CJK SC}
\setmonofont{Noto Sans Mono CJK SC}
\setCJKmainfont{Noto Sans CJK SC}
\title{BraTS MRI Reconstruction}
\subtitle{BME AI Project 1}
\author{第4组 / 唐志昊 2022533131}
\date{2026-04-28}

\begin{document}

\begin{frame}
  \titlepage
\end{frame}

\begin{frame}{Project Motivation}
\begin{itemize}
  \item Faster MRI acquisition is clinically desirable but undersampling causes aliasing artifacts.
  \item Goal: reconstruct high-quality T2 MRI from AF=5 undersampled k-space.
  \item Tasks: simulation, baseline reconstruction, and multi-modal unrolled reconstruction.
\end{itemize}
\end{frame}

\begin{frame}{Dataset and Experimental Setup}
\begin{itemize}
  \item Dataset: BraTS, using co-registered T1 and T2 modalities.
  \item Split: patient-level split to avoid train/test leakage.
  \item Final formal run: 16 central slices per patient, all available patients.
  \item Runtime optimization: slice caching, pinned memory, TF32, non-blocking transfer.
\end{itemize}
\end{frame}

\begin{frame}{Task 1: Undersampling Simulation}
  \centering
  \includegraphics[width=0.92\linewidth,height=0.78\textheight,keepaspectratio]{outputs\_formal/task1/undersampling\_visualization.png}
\end{frame}

\begin{frame}{Task 2: Baseline U-Net}
\begin{columns}[T]
\column{0.48\linewidth}
\begin{itemize}
  \item Input: undersampled T2 slice.
  \item Model: U-Net baseline.
  \item Loss: MSE.
  \item Final performance: PSNR 35.18, SSIM 0.9271.
\end{itemize}
\column{0.48\linewidth}
  \includegraphics[width=\linewidth]{outputs\_formal/task2/loss\_curves.png}
\end{columns}
\end{frame}

\begin{frame}{Task 2 Reconstruction Examples}
  \centering
  \includegraphics[width=0.92\linewidth,height=0.78\textheight,keepaspectratio]{outputs\_formal/task2/reconstruction\_results.png}
\end{frame}

\begin{frame}{Task 3: Multi-modal Unrolled Network}
\begin{columns}[T]
\column{0.5\linewidth}
\begin{itemize}
  \item Input: undersampled T2 + fully sampled T1.
  \item Architecture: 2-cascade unrolled U-Net with data consistency.
  \item Loss: hybrid L1/L2.
  \item Final performance: PSNR 38.12, SSIM 0.9660.
\end{itemize}
\column{0.45\linewidth}
  \includegraphics[width=\linewidth]{outputs\_formal/task3/loss\_curves.png}
\end{columns}
\end{frame}

\begin{frame}{Quantitative Comparison}
\begin{columns}[T]
\column{0.52\linewidth}
\begin{itemize}
  \item Task 3 improves over Task 2 by 2.93 dB PSNR.
  \item Task 3 improves over Task 2 by 0.0389 SSIM.
  \item Multi-modal guidance and data consistency both contribute measurable gains.
\end{itemize}
\column{0.43\linewidth}
  \includegraphics[width=\linewidth]{outputs\_formal/report\_assets/metric\_comparison.png}
\end{columns}
\end{frame}

\begin{frame}{Task 3 Best Reconstructions}
  \centering
  \includegraphics[width=0.92\linewidth,height=0.78\textheight,keepaspectratio]{outputs\_formal/task3/best\_reconstructions.png}
\end{frame}

\begin{frame}{Task 3 Error Analysis}
  \centering
  \includegraphics[width=0.92\linewidth,height=0.78\textheight,keepaspectratio]{outputs\_formal/task3/error\_analysis.png}
\end{frame}

\begin{frame}{Conclusions}
\begin{itemize}
  \item The baseline model already removes most aliasing artifacts effectively.
  \item The multi-modal unrolled model provides the best fidelity and visual quality.
  \item Remaining hard cases are mainly complex slices with challenging local structure.
  \item Future work: stronger edge-aware loss, deeper cascades, and targeted hard-case training.
\end{itemize}
\end{frame}

\begin{frame}{AI Usage Declaration}
\begin{itemize}
  \item AI tools were used only for debugging support, automation, and draft editing.
  \item All model choices, experiment validation, metric checks, and final submitted materials were manually reviewed and confirmed.
  \item This declaration is included to satisfy the assignment requirement.
\end{itemize}
\end{frame}

\begin{frame}{Thank You}
  \centering
  Questions are welcome.
\end{frame}

\end{document}
